% Generated draft: requires mathematical and source review.
\documentclass[11pt]{article}
\usepackage[margin=25mm]{geometry}
\usepackage{fontspec}
\setmainfont{Latin Modern Roman}
\usepackage{amsmath,amssymb,mathtools,unicode-math}
\setmathfont{Latin Modern Math}
\usepackage{enumitem,needspace,xurl,hyperref}
\setlength{\parindent}{0pt}
\setlength{\parskip}{5pt}
\setlength{\emergencystretch}{3em}
\newcommand{\N}{\mathbb{N}}
\newcommand{\Z}{\mathbb{Z}}
\newcommand{\Q}{\mathbb{Q}}
\newcommand{\R}{\mathbb{R}}
\newcommand{\C}{\mathbb{C}}
\newcommand{\F}{\mathbb{F}}
\newcommand{\A}{\mathbb{A}}
\newcommand{\PP}{\mathbb{P}}
\newcommand{\E}{\mathbb{E}}
\newcommand{\eps}{\varepsilon}
\newcommand{\dd}{\,\mathrm{d}}
\begin{document}
\section*{UnsolvedMath ID: 700676}
\textbf{Source problem code:} catalogue entry
\subsection{Definitions and mathematical statement}
Let \(G\) be a group with identity \(1\), let \(H<G\) be a proper subgroup, and let \(a\in H\) satisfy \(a^2\ne 1\). Write \(a^g=g^{-1}ag\), and let \(\langle S\rangle\) denote the subgroup generated by a subset \(S\subseteq G\).

The source refers to Problem 6.55 for the definitions of a Frobenius group, a Frobenius complement, and a Frobenius kernel. That definition is not included in the supplied record. These terms below therefore retain their source-intended meanings; in particular, no finiteness or splitting hypothesis is added. The exact conventions for possibly infinite Frobenius groups and their kernels require verification against Problem 6.55.

Assume that
\[
\forall g\in G\setminus H,\quad F_g=\langle a,a^g\rangle
\text{ is a Frobenius group admitting a complement }C_g
\text{ with }a\in C_g.
\]
The complement \(C_g\) is required to contain \(a\); it is not required by this hypothesis to equal \(\langle a\rangle\).

Define the family of source-intended Frobenius kernels
\[
\mathcal K_a(G)=\bigl\{K:\text{ there exists a subgroup }F\le G
\text{ which is a Frobenius group with complement }\langle a\rangle
\text{ and corresponding kernel }K\bigr\},
\]
and its set-theoretic union
\[
U_a(G)=\bigcup_{K\in\mathcal K_a(G)}K.
\]
The question is whether, for every \((G,H,a)\) satisfying the stated hypotheses,
\[
U_a(G)\le G.
\]
This union ranges over all such Frobenius subgroups of \(G\), not merely the groups \(F_g\). Its precise interpretation remains conditional on the missing definitions from Problem 6.55.

The supplied editors' comment, dated 2005, reports an affirmative answer under either of the following additional conditions:
\begin{enumerate}
\item The order of \(a\) is even.
\item The order of \(a\) is neither \(3\) nor \(5\), and \(\langle a,a^g\rangle\) is finite for every \(g\in G\setminus H\).
\end{enumerate}
These are supplied historical assertions, not independently verified results in this draft. The registry labels the general problem as open.
\subsection{Short English statement}
Given a proper subgroup \(H<G\) and \(a\in H\) with \(a^2\ne1\), assume that each \(\langle a,a^g\rangle\), for \(g\notin H\), is a Frobenius group with a complement containing \(a\). Must the union of the kernels of all Frobenius subgroups of \(G\) having complement exactly \(\langle a\rangle\) be a subgroup? The source reports two affirmative special cases, but its referenced definitions for Frobenius groups and kernels must be checked.
\subsection{Sources}
\begin{itemize}
\item Supplied local UnsolvedMath registry record, permanent ID \(700676\), exact source code \texttt{catalogue entry}, title \emph{Kourovka Notebook Problem 10.61}; upstream revision \texttt{d371ee917ac9647fbb06cd6fc529deec6f3dda82}. The supplied status is \texttt{open}; the record was updated on 2026-10-16.
\item \emph{The Kourovka Notebook (Unsolved Problems in Group Theory)}, Problem 10.61, proposed by A. I. Sozutov in the 10th issue (1986). The supplied background identifies the 21st issue edition (2026), \texttt{arXiv:1401.0300v48}. The collection description attributes its editing to E. I. Khukhro and V. D. Mazurov.
\item Problem 6.55 of \emph{The Kourovka Notebook}, cited by the supplied statement for definitions; its text is not supplied.
\item A. I. Sozutov, \emph{Algebra and Logic}, \textbf{34}, no. 5 (1995), 295--305; supplied as a related reference. No article title is supplied.
\item Editors' comment (2005): A. M. Popov and A. I. Sozutov, \emph{Algebra and Logic}, \textbf{44}, no. 1 (2005), 40--45, cited for the affirmative result when the order of \(a\) is even. No article title is supplied.
\item Editors' comment (2005): A. M. Popov, \emph{Algebra and Logic}, \textbf{43}, no. 2 (2004), 123--127, cited for the affirmative result when the order of \(a\) is neither \(3\) nor \(5\) and all groups \(\langle a,a^g\rangle\), with \(g\notin H\), are finite. No article title is supplied.
\item All references, attribution, and historical assessments above are transcribed from the supplied record and have not been independently verified.
\end{itemize}
\end{document}
